% ServoModule
\section{【基本】ServoModule — PWM出力パターン}

\begin{patternbox}{ServoModule で学べるパターン}
\begin{itemize}
    \item LEDC（PWM）の初期化と制御
    \item 「変化がある場合のみ更新」による無駄な出力の抑制
    \item currentAngle による出力状態のフィードバック
\end{itemize}
\end{patternbox}

\subsection{Config/Data構造体}

\begin{lstlisting}[style=cppstyle, caption={ServoModule.h — Config/Data}]
struct ServoConfig {
    uint8_t  pin;            // サーボ信号ピン
    uint8_t  pwmChannel;     // LEDCチャネル番号（0-15）
    uint16_t minPulseUs;     // 最小パルス幅 [us]（例: 500）
    uint16_t maxPulseUs;     // 最大パルス幅 [us]（例: 2500）
    uint8_t  defaultAngle;   // 初期角度 [度]（0-180）
};

struct ServoData {
    uint8_t targetAngle  = 90;  // 目標角度（ロジックフェーズが設定）
    uint8_t currentAngle = 90;  // 現在の角度（モジュールが更新）
};
\end{lstlisting}

\begin{pointbox}[targetとcurrentの分離]
\texttt{targetAngle}はロジックフェーズが書き込む「指示値」、
\texttt{currentAngle}はモジュールが更新する「実際の出力値」。
この分離により、ロジックフェーズが「今サーボがどの角度にいるか」を確認できる。
\end{pointbox}

\subsection{updateOutput()の実装}

\begin{lstlisting}[style=cppstyle, caption={ServoModule.cpp — updateOutput()}]
void ServoModule::updateOutput(SystemData& data) {
    uint8_t target = constrain(data.servo.targetAngle, 0, 180);

    // 変化がある場合のみPWM更新
    if (target != data.servo.currentAngle) {
        ledcWrite(_config.pwmChannel, _angleToDuty(target));
        data.servo.currentAngle = target;
    }
}
\end{lstlisting}

\begin{pointbox}[変化検出パターン]
毎ループPWMを書き込むのではなく、\texttt{target != current}の場合のみ更新する。
同じ値を毎回書き込んでも動作に問題はないが、
不要なハードウェアアクセスを避けるのは良い習慣。
特にI2C/SPI経由のデバイスでは通信コストの削減になる。
\end{pointbox}

\begin{pointbox}[constrainによる安全性]
\texttt{constrain()}でサーボの物理的な可動範囲に制限している。
ロジックフェーズが誤って範囲外の値を設定しても、モジュール側で安全に丸める。
入力バリデーションはモジュールの責務。
\end{pointbox}

% ─── API仕様 ───────────────────────────────
\subsection{API仕様}


\begin{apibox}{ServoModule --- サーボモーター制御モジュール}
\textbf{定義ファイル:} \texttt{lib/ServoModule/ServoModule.h / .cpp}\\
\textbf{分類:} 出力モジュール（\texttt{updateOutput()}のみオーバーライド）\\
\textbf{対象デバイス:} PWMサーボモーター（SG90等）

LEDCによるPWMパルス幅制御でサーボモーターの角度を設定するモジュール。
角度（0--180$^{\circ}$）からパルス幅（$\mu$s）への変換を内部で行う。
\end{apibox}

\subsubsection{機能概要}

\begin{itemize}
    \item 角度指定（0--180$^{\circ}$）によるサーボ制御
    \item パルス幅範囲のカスタマイズ（\texttt{minPulseUs} / \texttt{maxPulseUs}）
    \item 角度変化時のみPWM更新（不要な書き込みを抑制）
    \item 初期角度の設定（\texttt{defaultAngle}）
\end{itemize}

% --- Config ---
\subsubsection{Config構造体}

\begin{funcblock}{ServoConfig}
サーボのPWMパラメータを定義する。

\begin{longtable}{l l l p{4.5cm}}
\toprule
\textbf{メンバ} & \textbf{型} & \textbf{制約・既定値} & \textbf{説明} \\
\midrule
\endhead
\texttt{pin}          & \texttt{uint8\_t}  & ---       & サーボ信号出力GPIOピン番号 \\
\texttt{pwmChannel}   & \texttt{uint8\_t}  & 0--15     & LEDCチャネル番号。他モジュールと重複しないこと \\
\texttt{minPulseUs}   & \texttt{uint16\_t} & 例: 500   & 0$^{\circ}$ 時のパルス幅 [$\mu$s] \\
\texttt{maxPulseUs}   & \texttt{uint16\_t} & 例: 2500  & 180$^{\circ}$ 時のパルス幅 [$\mu$s] \\
\texttt{defaultAngle} & \texttt{uint8\_t}  & 例: 90    & 初期角度 [度]（0--180） \\
\bottomrule
\end{longtable}
\end{funcblock}

% --- Data ---
\subsubsection{Data構造体}

\begin{funcblock}{ServoData}
サーボの目標角度と現在角度を保持する。

\begin{longtable}{l l l p{5.0cm}}
\toprule
\textbf{メンバ} & \textbf{型} & \textbf{初期値} & \textbf{説明} \\
\midrule
\endhead
\texttt{targetAngle}  & \texttt{uint8\_t} & \texttt{90} & 目標角度 [度]（0--180）。ロジックフェーズで設定 \\
\texttt{currentAngle} & \texttt{uint8\_t} & \texttt{90} & 現在のPWM出力に対応する角度 [度]（0--180） \\
\bottomrule
\end{longtable}
\end{funcblock}

% --- メソッド ---
\subsubsection{メソッド}

\begin{funcblock}{ServoModule(const ServoConfig\& config)}
\begin{tabularx}{\textwidth}{l X}
\toprule
\textbf{項目} & \textbf{内容} \\
\midrule
引数   & \texttt{config} --- \texttt{ServoConfig}へのconst参照 \\
動作   & Configを\texttt{\_config}にコピー保持する \\
\bottomrule
\end{tabularx}
\end{funcblock}

\begin{funcblock}{bool init()}
\begin{tabularx}{\textwidth}{l X}
\toprule
\textbf{項目} & \textbf{内容} \\
\midrule
戻り値   & \texttt{bool} --- 常に\texttt{true} \\
動作     & LEDCチャネルを50Hz（サーボ標準）で初期設定し、\texttt{defaultAngle}に対応するパルス幅を出力する \\
\bottomrule
\end{tabularx}
\end{funcblock}

\begin{funcblock}{void updateOutput(SystemData\& data)}
\begin{tabularx}{\textwidth}{l X}
\toprule
\textbf{項目} & \textbf{内容} \\
\midrule
読み取り元     & \texttt{data.servo} \\
タイミング制御 & 毎ループ実行（ただし角度変化なしの場合はPWM更新をスキップ） \\
動作           & \texttt{targetAngle}と\texttt{currentAngle}を比較し、変化があれば角度をパルス幅に変換してPWMデューティを更新。\texttt{currentAngle}を更新 \\
\bottomrule
\end{tabularx}
\end{funcblock}

\begin{notebox}[角度→パルス幅変換]
角度からパルス幅への変換は以下の線形補間で行う:

\vspace{0.3em}
$$pulseWidth = minPulseUs + \frac{angle}{180} \times (maxPulseUs - minPulseUs)$$
\vspace{0.3em}

\noindent SG90の場合（500--2500$\mu$s）: 90$^{\circ}$ $\rightarrow$ 1500$\mu$s（中央）
\end{notebox}

% --- 使用例 ---
\subsubsection{使用例}

\begin{lstlisting}[style=cppstyle, caption={ServoModuleの使用例}]
// ProjectConfig.h
const ServoConfig SERVO_CONFIG = {
    .pin          = 1,     // GPIO1: サーボ信号
    .pwmChannel   = 0,     // LEDCチャネル0
    .minPulseUs   = 500,   // SG90: 500us
    .maxPulseUs   = 2500,  // SG90: 2500us
    .defaultAngle = 90,    // 初期角度 90度（中央）
};

// ロジックフェーズでの使用
void applyPattern(SystemData& data) {
    // IMUの傾きに応じてサーボ角度を設定
    if (data.mpu.isValid) {
        float tilt = atan2(data.mpu.accelX, data.mpu.accelZ);
        data.servo.targetAngle = constrain(
            map(tilt * 57.3, -45, 45, 0, 180), 0, 180
        );
    }
}
\end{lstlisting}
